%\chapter{Analysis of Strain} \setcounter{page}{32} \section{Extension}
Whenever, owing to any cause, changes take place in the relative positions 
of the parts of a body the body is said to be \textit{strained}. A very simple 
example of a strained body is a stretched bar. Consider a bar of square 
section suspended vertically and loaded with a weight at its lower end. 
Let a line be traced on the bar in the direction of its length, let two points of 
the line be marked, and let the distance between these points be measured. 
When the weight is attached the distance in question is a little greater than 
it was before the weight was attached. Let \(l_o\) be the length before stretching, and
\(l\) the length when stretched. Then \((l-l_o)/l_o\), is a number (generally a 
very small fraction) which is called the \textit{extension} of the line in question. If 
this number is the same for all lines parallel to the length of the bar, it may 
be called \textit{the extension of the bar}. A steel bar of sectional area \(1~in^2\)
\(=6.4515~cm^2\) loaded with 1 ton \(= 1016.05 kg\) will undergo an extension of about
\(7 \times 10^{-5}\). It is clear that for the measurement of such 
small quantities as this rather elaborate apparatus and refined methods of 
observation are required \footnote{See, for example, Ewing, Strength of
Materials, Cambridge, 1899, pp. 73 et seq., or G. P. C. Searle, Experimental
Elasticity, Cambridge, 1908.} Without attending to methods of measurement 
we may consider a little more in detail the state of strain in the stretched 
bar. Let \(e\) denote the extension of the bar, so that its length is increased in 
the ratio \(1+e:1\), and consider the volume of the portion of the bar
contained between any two marked sections. This volume is increased by
stretching the bar, but not in the ratio \(1+e:1\) . When the bar is stretched
longitudinally, it contracts laterally. If the linear lateral contraction is
\(e^{\prime}\), the sectional area is diminished in the ratio
\((1-e^{\prime})^2:1\), and the volume in question is increased in the ratio
\((1+e)(1-e^{\prime})^2:1\).  In the case of a bar under tension,
\(e^{\prime}\) is a certain multiple of \(e\), say \(\sigma e\)
\footnote{Kirchhoff's notation, Poisson's ratio, Filon used \(\nu\).},
and \(\sigma\) is about \(1/3\) for very many materials. If \(e\)
is very small and \(e^2\) is neglected, the areal contraction \(2 \sigma e\), 
\[
(1-2 \sigma e)(1+e) = 1 + e - 2 \sigma e  = 1+ e(1-2 \sigma),
\]
and, the cubical dilatation is \((1-2 \sigma) e\).

For the analytical description of the state of strain in the bar we should 
take an origin of coordinates \(x\), \(y\), \(z\) on the axis, and measure the
coordinate \(z\)

along the length of the bar. Any particle of the bar which has the coordinates
\(x\), \(y\), \(z\)
when the weight is not attached will move after the attachment 
of the weight into a new position. Let the particle which was at the origin 
move through a distance
\(z_o\), then the particle which was at
\((x,y,z)\)
moves to 
the point of which the coordinates are 
\[
(x(1-\sigma e), \quad y(1-\sigma e), \quad z_o + (z-z_o) (1+ e)).
\]
The state of strain is not very simple. If lateral forces could be applied to 
the bar to prevent the lateral contraction the state of strain would be very 
much simplified. It would then be described as a \textit{simple extension}. 

\section{Pure  shear}
As a second example of strain let us suppose that lateral forces are applied 
to the bar so as to produce extension of amount
\(\epsilon_1\) of lines parallel to the axis of \(x\) and
extension of amount \(\epsilon_2\) of lines parallel to the axis \(y\), and that 
longitudinal forces are applied, if any are required, to prevent any extension 
or contraction parallel to the axis of \(z\). The particle which was at
\((x, y, z)\) will move to \((x+\epsilon_1 x,y+\epsilon_2 y, z)\)
and the area of the section will be increased in the ratio
\((1+\epsilon_1)(1+\epsilon_2):1\).
If \(\epsilon_1\) and \(\epsilon_2\)
are related so that this ratio is equal to unity,
\[
(1+\epsilon_1)(1+\epsilon_2)=1,
\]
there will be no change in the area of the section or in the volume of any
portion of the bar, but the shape of the section will be distorted.  Either
\(\epsilon_1\) or \(\epsilon_2\) is then negative, or there is contraction of
the corresponding set of lines.  The strain set up in the bar is called
\textit{pure shear}.  Fig. 1 below shows a square \(ABCB\) distorted by pure
shear into a rhombus \(A'B'C'D'\) of the same area. 

\section{Simple shear}
As a third example of strain let us suppose that the bar after being distorted
by pure shear is turned bodily about its axis. We suppose that the axis of
\(x\) is the direction in which contraction takes place, and we put 
\[
\epsilon_2 - \epsilon_1 = 2 \tan \alpha.
\]

Then we can show that, if the rotation is of amount \(\alpha\) in the sense
from \(y\) to \(x\), the position reached by any particle is one that could
have been reached by the sliding of all the particles in the direction of a
certain line through distances proportional to the distances of the particles
from a certain plane containing this line. 

Since \((1+\epsilon_1)(1+\epsilon_2)=1\), and \(\epsilon_2 -\epsilon_1 = 2 \tan
\alpha\), we have 
\[
1 + \epsilon_1 = \sec(\alpha ) - \tan(\alpha ), \quad
1 + \epsilon_2 = \sec(\alpha ) + \tan(\alpha ).
\] 
By the pure shear, the particle which was at
\((x, y)\)
will move to
\((x_1,y_1)\), where 
\[
\begin{bmatrix}
x_1\\
y_1 
\end{bmatrix}=
\begin{bmatrix}
(\sec(\alpha ) - \tan(\alpha ))x\\
(\sec(\alpha ) + \tan(\alpha ))y 
\end{bmatrix}
\] 
and by the rotation it is moved again to
\((x_2,y_2)\), where
\[
\begin{bmatrix} x_2\\ y_2 \end{bmatrix}=
\begin{bmatrix}
\cos(\alpha ) & \sin(\alpha )\\
-\sin(\alpha ) & \cos(\alpha )
\end{bmatrix}
\begin{bmatrix} x_1\\ y_1 \end{bmatrix}.
\] 
so that we have
\[
\begin{bmatrix} x_2\\ y_2 \end{bmatrix}=
\begin{bmatrix} x\\ y \end{bmatrix}+ \tan(\alpha) 
\begin{bmatrix}
-\cos(\alpha )  & 1+\sin(\alpha )\\
-1+\sin(\alpha )& \cos(\alpha )
\end{bmatrix}
\begin{bmatrix}
x\\
y
\end{bmatrix}.
\]

Now, writing 
\(\beta\) for 
\(\frac{1}{2}\pi-a\), we have 
\[
\begin{bmatrix} x_2\\ y_2 \end{bmatrix}=
\begin{bmatrix} x\\ y \end{bmatrix}+ 
2 \tan(\alpha) \left(
-x \sin \frac{\beta}{2} + y \cos\frac{\beta}{2}
\right)
\begin{bmatrix}
\cos(\beta/2) \\
\sin(\beta/2) 
\end{bmatrix}
\]
and we can observe that 
\[
\begin{aligned}
x_2 \cos \frac{\beta}{2} + y_2 \sin \frac{\beta}{2}=&
x \cos \frac{\beta}{2} + y \sin \frac{\beta}{2}
+ 2 \tan(\alpha) \left(
-x \sin \frac{\beta}{2} + y \cos\frac{\beta}{2} \right), \\
-x_2 \sin \frac{\beta}{2} + y_2 \cos\frac{\beta}{2}&=
-x\sin\frac{1}{2}\beta+y\cos\frac{1}{2}\beta.
\end{aligned}
\]
Hence, taking axes of \(X\) and \(Y\) which are obtained from those of
\(x\) and \(y\) by a rotation through 
\(\frac{\pi}{4}\pi - \frac{\alpha}{2}\) in the sense from
\(x\) towards \(y\), we see that the particle which was at
\((X, Y)\) is moved by the pure shear followed by 
the rotation to the point
\((X_2, Y_2)\), where
\[
\begin{bmatrix}
X_2 \\
Y_2
\end{bmatrix} =
\begin{bmatrix}
1 & 2\tan\alpha \\
0 &  1
\end{bmatrix}
\begin{bmatrix}
X \\
Y
\end{bmatrix}
.\]

Thus every plane of the material which is parallel to the plane of
\((X, z)\) slides 
along itself in the direction of the axis of
\(X\) through a distance proportional 
to the distance of the plane from the plane of
\((X, z)\). The kind of strain just 
described is called a \textit{simple shear}, the angle \(\alpha\)
is the \textit{angle of the shear}, 
and \(2 \tan \alpha\) is the \textit{amount of the shear}.

\begin{figure}[!h]
\caption{Shear}
 \label{fig:1}
\end{figure}
Fig. 1 shows a square \(ABCB\) distorted by pure shear into a rhombus
\(A'B'C'D'\) of the same area, which is then rotated into the position
\(A''B''C''D''\). The angle of the shear is \(A'OA''\), and the angle \(AOX\)
is half the complement of this angle.  The lines \(AA'',BB'',CC'',DD''\) are
parallel to \(OX\) and proportional to their distances from it.

We shall find that all kinds of strain can be described in terms of
simple extension and simple shear, but for the discussion of complex
states of strain and for the expression of them by means of simpler
strains we require a general kinematical theory\footnote{The greater
  part of the theory is due to Cauchy (see Introduction). Some
  improvements were made by Clebsch in his treatise of 1862, and others
  were made by Kelvin and Tait, Nat. Phil. Part I.}

\section{Displacement.} % 4

We have, in every case, to distinguish two states of a body a first
state and a second state. The particles of the body pass from their
positions in the first state to their positions in the second state by a
displacement. The displacement may be such that the line joining any two
particles of the body has the same length in the second state as it has
in the first; the displacement is then one which would be possible in a
rigid body. If the displacement alters the length of any line, the
second state of the body is described as a \textit{strained state}, and
then the first state is described as the \textit{unstrained state}.

In what, follows we shall denote the coordinates of the point occupied
by a particle, in the unstrained state of the body, by x, y, z, and the
coordinates of the point occupied by the same particle in the strained
state by \(x + u, y+v, z + w.\) Then \(u, v, w\) are the projections on
the axes of a vector quantity --- the displacement. We must take
\(u, v, w\) to be continuous functions of \(x, y, z,\) and we shall in
general assume that they are analytic functions.

It is clear that, if the displacement (u, v, w) is given, the strained
state is entirely determined; in particular, the length of the line
joining any two particles can be determined.

\section{Displacement in simple extension and simple shear.} % 5
The displacement in a simple extension parallel to the axis of x is
given by the equations \(u=ex, v=0, w=0,\) where e is the amount of
the extension. If e is negative there is contraction.

The displacement in a simple shear of amount (=2 tan a), by which lines
parallel to the axis of x slide along themselves, and particles in any
plane parallel to the plane of (X, y) remain in that plane, is given by
the equations \(u=s y, v=0, w=0.\)

In Fig. 2, \(A B\) is a segment of a line parallel to the axis of \(x\), which
subtends an angle \(2 a\) at \(O\) and is bisected by \(O y\). By the simple
shear particles lying on the line \(O A\) are displaced so as to lie on \(O
B\). The particle at any point \(P\) on \(A B\) is displaced to \(Q\) on \(A
B\) so that \(P Q=A B\), and the particles on \(O P\) are displaced to points
on \(O Q\). A parallelogram such as \(O P N M\) becomes a parallelogram such as
\(O Q K M\).

